\documentclass[11pt]{article}

\usepackage[margin=1in]{geometry}
\usepackage[T1]{fontenc}
\usepackage[utf8]{inputenc}
\usepackage{lmodern}
\usepackage{xcolor}
\usepackage{booktabs}
\usepackage{longtable}
\usepackage{array}
\usepackage{amsmath,amssymb}
\usepackage{enumitem}
\usepackage{hyperref}

\hypersetup{colorlinks=true, linkcolor=blue, urlcolor=blue}
\definecolor{accent}{RGB}{25,56,99}

\setlength{\parskip}{0.45em}
\setlength{\parindent}{0pt}

\begin{document}

\begin{center}

{\LARGE\bfseries\color{accent}
Comparison of Power-Law and Logistic Models for Conductivity–Saturation Relationships}\\[0.4em]

{\large UNORM analysis of SIP column experiments at 0.01 Hz}

\end{center}

\section{Objective}

This report evaluates two competing functional models describing the relationship between water saturation ($S$) and electrical conductivity ($\sigma$) in Spectral Induced Polarization (SIP) column experiments.

The analysis focuses exclusively on the \textbf{UNORM conductivity relationship}:

\[
\sigma \quad \text{vs} \quad S
\]

using the dataset extracted from the SIP processing pipeline.

Two candidate models are compared:

\begin{itemize}

\item Power-law model
\item Logistic (sigmoidal) model

\end{itemize}

Model performance is evaluated using likelihood-based statistics and information criteria in order to determine which functional representation best explains the observed behavior.

\section{Data sources}

The dataset originates from processed SIP experiment folders generated by the SIP pipeline. Each run provides measurements of:

\begin{itemize}

\item Water saturation ($S$)
\item Electrical conductivity ($\sigma$)
\item Associated metadata (run name, site name)

\end{itemize}

The comparison analysis is performed at three aggregation levels:

\begin{longtable}{p{0.20\textwidth} p{0.65\textwidth}}
\toprule
\textbf{Level} & \textbf{Description}\\
\midrule
\endhead

Sample & Each experimental run is analyzed independently.\\

Site & All runs from the same site are combined into a single dataset.\\

All & All available runs are pooled together into a global dataset.\\

\bottomrule
\end{longtable}

This hierarchy allows identification of both local and global saturation–conductivity trends.

\section{Model formulations}

\subsection{Power-law model}

The power-law model is defined as:

\[
\sigma = a S^{b}
\]

where

\begin{itemize}

\item $a$ is a scaling constant
\item $b$ is the saturation exponent

\end{itemize}

Taking logarithms yields a linear regression form:

\[
\log \sigma = b \log S + \log a
\]

The workflow supports either base-10 logarithms or natural logarithms.

\subsection{Logistic model}

The logistic model captures sigmoidal behavior:

\[
\sigma =
c + \frac{L}{1+\exp[-k(S-x_0)]}
\]

where

\begin{itemize}

\item $L$ controls amplitude
\item $k$ controls steepness
\item $x_0$ defines the transition saturation
\item $c$ defines the baseline conductivity

\end{itemize}

This model can represent threshold-driven electrical conduction where conductivity increases rapidly after a critical saturation level.

Unlike earlier implementations, the logistic amplitude is allowed to be positive or negative, enabling the model to represent both increasing and decreasing relationships.

\section{Statistical evaluation}

Several statistical metrics are used to evaluate model performance.

\subsection{Sum of squared errors}

\[
SSE = \sum (y_i - \hat{y}_i)^2
\]

Lower SSE values indicate better agreement between predicted and observed values.

\subsection{Residual standard deviation}

\[
\sigma = \sqrt{\frac{SSE}{n}}
\]

This quantity represents the typical magnitude of residual errors.

\subsection{Gaussian log-likelihood}

Assuming normally distributed residuals:

\[
\ln L =
-\frac{n}{2}
\left[
\ln(2\pi\sigma^2)+1
\right]
\]

Higher log-likelihood values indicate a better statistical fit.

\subsection{Akaike Information Criterion}

Model comparison is primarily based on the Akaike Information Criterion:

\[
AIC = n\ln\left(\frac{SSE}{n}\right) + 2k
\]

where $k$ is the number of model parameters.

\begin{itemize}

\item Power-law model: $k=2$
\item Logistic model: $k=4$

\end{itemize}

The preferred model is the one with the \textbf{smaller AIC value}.

\section{Normalization strategy}

Normalization may optionally be applied before regression. When enabled, normalization is performed only within the relevant aggregation level:

\begin{itemize}

\item sample-level normalization uses sample data
\item site-level normalization uses site data
\item global normalization uses the full dataset

\end{itemize}

This prevents scale distortions when comparing models across different aggregation levels.

Two normalization methods are available:

\subsection{Z-score normalization}

\[
x' = \frac{x-\mu}{\sigma}
\]

\subsection{Min–max normalization}

\[
x' = \frac{x-x_{\min}}{x_{\max}-x_{\min}}
\]

\section{Interpretation of conductivity behavior}

Conductivity in partially saturated porous media is often influenced by pore connectivity.

At low saturation levels, water films may be disconnected, limiting conductive pathways. As saturation increases, continuous pathways develop and conductivity can increase rapidly.

This behavior produces a sigmoidal relationship between saturation and conductivity that may be better represented by logistic models.

In contrast, power-law relationships imply continuous scaling without a clear transition threshold.

\section{Outputs}

The comparison workflow generates several outputs:

\begin{longtable}{p{0.40\textwidth} p{0.18\textwidth} p{0.30\textwidth}}
\toprule
\textbf{File} & \textbf{Scope} & \textbf{Description}\\
\midrule
\endhead

\texttt{compare\_cond\_unorm\_\_0p01Hz.csv} & Global & Table containing fitted parameters and statistics.\\

\texttt{compare\_cond\_unorm\_\_0p01Hz.xlsx} & Global & Excel workbook summarizing all comparisons.\\

\texttt{compare\_cond\_unorm\_used\_data\_\_0p01Hz.csv} & Global & Dataset used for model fitting.\\

Plots & Per aggregation level & Visual comparison of conductivity vs saturation fits.\\

\bottomrule
\end{longtable}

Each row in the output tables includes:

\begin{itemize}

\item fitted model parameters
\item residual standard deviation
\item log-likelihood
\item AIC values
\item preferred model

\end{itemize}

\section{Summary}

This report compares power-law and logistic representations of conductivity–saturation relationships using SIP column experiment data.

The workflow:

\begin{itemize}

\item evaluates models at sample, site, and global scales
\item computes likelihood-based statistics
\item uses AIC for model selection
\item produces both tabular and graphical summaries
\end{itemize}

The results help identify whether conductivity follows smooth scaling behavior or exhibits threshold-driven transitions related to pore connectivity.

\end{document}
